\wysubtitle{logos.pdf}{出题：\textsf{XXY \& CH \& ZCH \& QHY}}
{https://qm.qq.com/q/OlkSiy9W04}{HDU 数学营}
{http://weixin.qq.com/r/mhKYgFTEMQlOrRDj90eI}{未央学社公众号}
\begin{flushleft}
    \kaishu\qquad 本资料仅作为模拟练习之用，目的是为了帮助大家更有效地复习，并减轻对考试的担忧。请正确的对待此资料，其旨在辅助复习，而非预示具体的考试内容。我们鼓励同学们认真复习，大学学习主打理解，而非刷题，期望大家在期中考试中取得优异成绩。
\end{flushleft}

\section{填空题}
\begin{problem}[行列式及其运算]
    设$\mat{A}$为四阶方阵，且$\abs{\mat{A}}=-3$，则$\abs{2\mat{A}^*}=$\ans{$-432$}.
\end{problem}
\solute0
\begin{solution}*
    $\mat{A}^*=\abs{\mat{A}}\mat{A}^{-1}$，由此得$\abs{2\mat{A}^*}=\abs{2\cdot(-3)\mat{A}^{-1}}=(-6)^4/\abs{\mat{A}}=-432$.
\end{solution}

\begin{problem}[特征值及特征向量]
    三阶方阵$\mat{A}$有特征值$\lambda_1=1,\ \lambda_2=2,\ \lambda_3=3$，则$\abs{\mat{A}^2-2\mat{A}+4\mat{E}}=$\ans{$84$}.
\end{problem}
\solute0
\begin{solution}*
    显然$\mat{A}^2-2\mat{A}+4\mat{E}$有特征值$3,\ 4,\ 7$，那么$\abs{\mat{A}^2-2\mat{A}+4\mat{E}}=3\times 4\times 7=84$.
\end{solution}

\begin{problem}[相似矩阵]
    已知矩阵$\mat{A=PQ}$，其中$\mat{P}=(1,2,1)^\T,\ 
    \mat{Q}=(2,-1,2)$. 则$\mat{A}^{100}=$\ans{$2^{99}\begin{pmatrix}
        2 & -1 & 2\\4 & -2 & 4\\2 & -1 & 2
    \end{pmatrix}$}.
\end{problem}
\solute1
\begin{solution}*
    $\mat{A}^{100}=\mat{PQPQ}\cdots\mat{QPQ}=\mat{P(QPQP}\cdots)\mat{Q}=\mat{P}\cdot 2^{99}\mat{Q}=2^{99}\begin{pmatrix}
        2 & -1 & 2\\4 & -2 & 4\\ 2 & -1 & 2
    \end{pmatrix}$.
\end{solution}

\begin{problem}[向量组及其秩]
    若矩阵$\mat{A}=\begin{pmatrix}
        2 & 1 & 0\\1 & 2 & 0\\0 & 1 &x
    \end{pmatrix}$
    与$\mat{B}=\begin{pmatrix}
        1 & 2 & -1\\2 & 1 & 3\\3 & 3 & 1
    \end{pmatrix}$等价，则$x=$\ans{$0$}.
\end{problem}
\solute0
\begin{solution}*
    易知$\abs{\mat{B}}=0$，所以$\abs{\mat{A}}=0$，故$x=0$.
\end{solution}

\begin{problem}[线性方程组的解释理论]
    设$\mat{A}$是$5\times 4$矩阵，$\rank{\mat{A}}=2$，$\eta_1,\ \eta_2,\ \eta_3$为非齐次线性方程组$\mat{A}\bm{x}=\bm{b}$的三个相异解且$\eta_1-\eta_2,\ \eta_2-\eta_3$相互独立，则$\mat{A}\bm{x}=\bm{b}$的通解为\ans{$\eta_1+\lambda_1(\eta_1-\eta_2)+\lambda_2(\eta_2-\eta_3)$}.
\end{problem}
\begin{solution}*
    $\rank{\mat{A}}=2,\ n=4$，可知$\mat{A}\bm{x}=\bf{0}$的解空间秩为$n-\rank{\mat{A}}=2$. 显然$\eta_1-\eta_2,\ \eta_2-\eta_3$为$\mat{A}\bm{x}=\bf{0}$的两个相互独立的解. 故$\mat{A}\bm{x}=\bm{b}$的通解为$\eta_1+\lambda_1(\eta_1-\eta_2)+\lambda_2(\eta_2-\eta_3)$.
\end{solution}


\begin{problem}
    设$\mat{A},\mat{B}$均为$n$阶正交矩阵，若$\abs{\mat{A}}=1,\ \abs{\mat{B}}=-1$，则$\abs{\mat{A+B}}=$\ans{$0$}.
\end{problem}
\solute0
\begin{solution}*
    $\abs{\mat{A+B}}=\abs{\mat{A}}\abs{\mat{E+A^\T B}}=\abs{\mat{A}}\abs{\mat{E+B^\T A}}=-\abs{\mat{B}}\abs{\mat{E+B^\T A}}=-\abs{\mat{A+B}}$. 所以$\abs{\mat{A+B}}=0$.
\end{solution}

\section{选择题}
\begin{choice}{D}
    若$\mathrm{D}=\begin{vmatrix}
        a_{11} & a_{12} & a_{13}\\
        a_{21} & a_{22} & a_{23}\\
        a_{31} & a_{32} & a_{33}
    \end{vmatrix}$，则行列式
    $\mathrm{B}=\begin{vmatrix}
         a_{11} & 2a_{12} & 3a_{13}\\
        2a_{21} & 4a_{22} & 6a_{23}\\
        3a_{31} & 6a_{32} & 9a_{33}
    \end{vmatrix}$的值为
\begin{tasks}(4)
    \task $6\mathrm{D}$
    \task $12\mathrm{D}$
    \task $24\mathrm{D}$
    \task $36\mathrm{D}$
\end{tasks}
\end{choice}
\solute0
\begin{solution}
    将$\mat{D}$的第$2$行及第$2$列$\times 2$，第$3$行及第$3$列$\times 3$则得到$\mat{B}$. 所以$\abs{\mat{B}}=\abs{\mat{D}}\cdot 2^2\cdot 3^2=36\mathrm{D}$. \sokka{D}
\end{solution}

\begin{choice}{A}
    设$\bm{\alpha}$为$n$维单位列向量，$E$为$n$阶单位矩阵，则
\begin{tasks}(4)
    \task $\mat{E}-\bm{\alpha\alpha}^\T$不可逆
    \task $\mat{E}+\bm{\alpha\alpha}^\T$不可逆
    \task $\mat{E}-2\bm{\alpha\alpha}^\T$不可逆
    \task $\mat{E}+2\bm{\alpha\alpha}^\T$不可逆
\end{tasks}
\end{choice}
\begin{solution}
    显然若$\bm{\alpha}_i=\begin{cases}
        1, & i=m\\0, & i\neq m
    \end{cases}$，则$\bm{\alpha\alpha}^\T_{ij}=\begin{cases}
        1, & i=j=m\\0, & i\neq m \text{ and } j\neq m
    \end{cases}$. 易知$\abs{\mat{E}-\bm{\alpha\alpha}^\T}=0$. \sokka{A}
\end{solution}

\begin{choice}{D}
    如果$\begin{cases}
        kx+z=0\\
        2x+ky+z=0\\
        kx-2y+z=0
    \end{cases}$
    有非零解，则$k$应为
\begin{tasks}(4)
    \task $-2$
    \task $-1$
    \task $1$
    \task $2$
\end{tasks}
\end{choice}
\solute0
\begin{solution}*
    系数矩阵行列式的值为$4-2k$，若$k=2$则行列式为$0$，系数矩阵不满秩，有非零解. \sokka{D}
\end{solution}

\begin{choice}{D}
    设$\mat{A}$为$n$阶矩阵，$\bm{\alpha}$为$n$维列向量. 若$\mat{B}=\begin{pmatrix}
        \mat{A} & \bm{\alpha}\\\bm{\alpha}^\T & \mat{O}
    \end{pmatrix},\ \rank{\mat{B}}=\rank{\mat{A}}$，则线性方程组
\begin{tasks}(2)
    \task $\mat{A}\bm{x}=\bm{\alpha}$必有无穷多解
    \task $\mat{A}\bm{x}=\bm{\alpha}$必有唯一解
    \task $\mat{B}(\mat{X},\mat{Y})^\T=\mat{O}$仅有零解
    \task $\mat{B}(\mat{X},\mat{Y})^\T=\mat{O}$必有非零解
\end{tasks}
\end{choice}
\begin{solution}*
    由于$\rank{\mat{B}}<n+1$，所以$\mat{B}(\mat{X},\mat{Y})^\T=\mat{O}$必有非零解.
\end{solution}

\begin{choice}{B}
    设$n$维向量组$\bm{s}_1,\ \bm{s}_2,\ \bm{s}_1\cup\bm{s}_2$的秩为$r,\ s,\ t$，则有
\begin{tasks}(4)
    \task $t\geq r+s$
    \task $t\leq r+s$
    \task $t=r+s$
    \task 不确定
\end{tasks}
\end{choice}
\begin{solution}*
    设$\bm{s}_1$的最大无关组是$\bm{\xi}_1,\ \bm{\xi}_2,\ \cdots,\ \bm{\xi}_r$，$\bm{s}_2$的最大无关组是$\bm{\eta}_1,\ \bm{\eta}_2,\ \cdots,\ \bm{\eta}_s$，显然对于任意向量$\bm{u}\in\bm{s}_1\cup\bm{s_2}$，$\bm{u}$都能被以上向量线性表示，所以$\mathrm{R}(\bm{s}_1\cup\bm{s}_2)=t\leq r+s$.
\end{solution}

\begin{choice}{A}
    设$\mat{A},\ \mat{B}$均为$n$阶实对称矩阵，则以下命题正确的是
\begin{tasks}[label=\roman*.](2)
    \task $\mat{A}$与$\mat{B}$相似$\Longleftrightarrow\abs{\lambda\mat{E}_n-\mat{A}}=\abs{\lambda\mat{E}_n-\mat{B}}$
    \task 若$\mat{A},\ \mat{B}$合同，则$\mat{A}$与$\mat{B}$相似
    \task! 若$\mat{A},\ \mat{B}$相似，则$\mat{A}$与$\mat{B}$合同
    \task! $n$元二次型$\bm{x}^\T\mat{A}\bm{x}$与$\bm{x}^\T\mat{B}\bm{x}$具相同规范型$\Longleftrightarrow\abs{\lambda\mat{E}_n-\mat{A}}=\abs{\lambda\mat{E}_n-\mat{B}}$
\end{tasks}
\begin{tasks}(4)
    \task i, ii
    \task ii, iii
    \task i, iii
    \task i, iv
\end{tasks}
\end{choice}
\begin{solution}
\begin{tasks}[label=\roman*.](2)
    \task 显然正确.
    \task 显然错误，例如$\mat{A}$与$2\mat{A}$具有相同规范型.
    \task! $\mat{PAP}^{-1}=\mat{B}$，若$\mat{A},\ \mat{B}$实对称，则存在$\mat{U},\ \mat{UU}^\T=\mat{E}$，$\mat{U}^\T\mat{AU}=\Lambda_1,\ \mat{V}^\T\mat{BV}=\Lambda_2$. 若$\mat{A},\ \mat{B}$相似，则$\Lambda_1=\Lambda_2$，所以$\mat{UAU}^\T=\mat{V}^\T\mat{BV},\ \mat{B}=\mat{VUAU}^\T\mat{V}^\T$，故$\mat{A},\ \mat{B}$合同.
    \task! 若$\mat{B}=\mat{V}^\T\mat{A}\mat{V}$而$\mat{V}$不正交，则$\mat{A},\ \mat{B}$不相似. \sokka{A}
\end{tasks}
\end{solution}

\section{计算题}
\begin{problem}
    设$\mat{A}=\begin{pmatrix}
        1 & 1 & -1\\
        -1 & 1 & 1\\
        1 & -1 & 1
    \end{pmatrix}$
    且$\mat{A}^*\mat{X}=\mat{A}^{-1}+2\mat{X}$，求矩阵$\mat{X}$.
\end{problem}
\solute1
\begin{solution}*
    $\abs{A}=4$，所以$\mat{A}$可逆，$\mat{A}^*=\abs{\mat{A}}\mat{A}^{-1}=4\mat{A}^{-1},\ \mat{X}=\ab[\mat{A}(4\mat{A}^{-1}-2\mat{E})]^{-1}=(4\mat{E}-2\mat{A})^{-1}$.
    \[(4\mat{E}-2\mat{A}:\mat{E})=\begin{pNiceArray}{ccc:ccc}
        2 & -2 & 2 & 1 & 0 & 0\\
        2 & 2 & -2 & 0 & 1 & 0\\
        -2 & 2 & 2 & 0 & 0 & 1
    \end{pNiceArray}\overset{r}{\sim}
    \begin{pNiceArray}{ccc:ccc}
        1 & 0 & 0 & \tfrac14 & \tfrac14 & 0\\
        0 & 1 & 0 & 0 & \tfrac14 & \tfrac14\\
        0 & 0 & 1 & \tfrac14 & 0 & \tfrac14
    \end{pNiceArray},\ \mat{X}=\begin{pmatrix}
        \tfrac14 & \tfrac14 & 0\\
        0 & \tfrac14 & \tfrac14\\
        \tfrac14 & 0 & \tfrac14
    \end{pmatrix}\]
\end{solution}
\begin{problem}
    已知向量组$\bm{\alpha}_1=(1,1,1,3)^\T,\ \bm{\alpha}_2=(1,3,-5,-1)^\T,\ \bm{\alpha}_3=(3,1,10,15)^\T,\ \bm{\alpha}_4=(3,7,-9,3-2a)^\T$线性相关，求出$a$的值，并写出$\bm{\alpha}_4$由$\bm{\alpha}_1,\ \bm{\alpha}_2,\ \bm{\alpha}_3$的表达式.
\end{problem}
\solute1
\begin{solution}
    \[\begin{pmatrix}
        1 & 1 & 3 & 3\\
        1 & 3 & 1 & 7\\
        1 & -5 & 10 & -9\\
        3 & -1 & 15 & 3-2a
    \end{pmatrix}\sim
    \begin{pmatrix}
        1 & 0 & 0 & 1\\
        0 & 1 & 0 & 2\\
        0 & 0 & 1 & 0\\
        0 & 0 & 0 & 2-2a
    \end{pmatrix}\]
    所以$a=1,\ (1,2,0,0)^\T=(1,0,0,0)^\T+2(0,1,0,0)^\T$，得$\bm{\alpha}_4$的表达式为$\bm{\alpha}_4=\bm{\alpha}_1+2\bm{\alpha}_2$.
\end{solution}

\begin{problem}
    已知三维向量空间中有两个基$\bm{\alpha}_1=(1,1,1)^\T,\ \bm{\alpha}_2=(1,0,-1)^\T,\ \bm{\alpha}_3=(1,0,1)^\T$与$\bm{\beta}_1=(1,2,1)^\T,\ \bm{\beta}_2=(2,3,4)^\T,\ \bm{\beta}_3=(3,4,3)^\T$.
\begin{enumerate}
    \item 求基$\bm{\alpha}_1,\ \bm{\alpha}_2,\ \bm{\alpha}_3$到基$\bm{\beta}_1,\ \bm{\beta}_2,\ \bm{\beta}_3$的过渡矩阵$\mat{P}$.
    \item 设$\xi$在基$\bm{\alpha}_1,\ \bm{\alpha}_2,\ \bm{\alpha}_3$的坐标为$(1,1,-1)^\T$，求$\xi$在基$\bm{\beta}_1,\ \bm{\beta}_2,\ \bm{\beta}_3$的坐标.
\end{enumerate}
\end{problem}
\solute1
\begin{solution}
\begin{enumerate}
    \item 令$\mat{A}=(\bm{\alpha}_1,\bm{\alpha}_2,\bm{\alpha}_3),\ \mat{B}=(\bm{\beta}_1,\bm{\beta}_2,\bm{\beta}_3)$，$\mat{AP}=\mat{B}$，则
    \[(\mat{A}:\mat{B})=\begin{pNiceArray}{ccc:ccc}
        1 & 1 & 1 & 1 & 2 & 3\\
        1 & 0 & 0 & 2 & 3 & 4\\
        1 & -1 & 1 & 1 & 4 & 3
    \end{pNiceArray}\to
    \begin{pNiceArray}{ccc:ccc}
        1 & 0 & 0 & 2 & 3 & 4\\
        0 & 1 & 0 & 0 & -1 & 0\\
        0 & 0 & 1 & -1 & 0 & -1
    \end{pNiceArray},\ \mat{P}=\begin{pmatrix}
        2 & 3 & 4\\
        0 & -1 & 0\\
        -1 & 0 & -1
    \end{pmatrix}\]
    \item 其中$\mat{P}^{-1}=\begin{pmatrix}
        -\tfrac12 & -\tfrac32 & -2\\
        0 & -1 & 0\\
        \tfrac12 & \tfrac32 & 1
    \end{pmatrix}$. 则$\xi$在基$\bm{\beta}_1,\ \bm{\beta}_2,\ \bm{\beta}_3$下坐标为$(y_1,y_2,y_3)^\T=\mat{P}^{-1}(x_1,x_2,x_3)^\T=(0,-1,1)^\T$.
\end{enumerate}
\end{solution}

\begin{problem}
    已知线性方程组$\begin{cases}
        x_1+x_2+x_3+x_4=-1\\
        4x_1+3x_2+5x_3-x_4=-1\\
        ax_1+x_2+3x_3+bx_4=1
    \end{cases}$有三个线性无关解.
\begin{enumerate}
    \item 求出方程组系数矩阵的秩.
    \item 求出$a,\ b$及方程组通解.
\end{enumerate}
\end{problem}
\begin{solution}
\begin{enumerate}
    \item 有三个线性无关解则解空间维数不小于$2$，所以$4-\rank{\mat{A}}\geq 2$. $\mat{A}$的子式$\begin{vmatrix}
        1 & 1\\4 & 3
    \end{vmatrix}\neq 0$，所以$\rank{\mat{A}}\geq 2$. 故$\rank{\mat{A}}=2$.
    \item 
    \[\begin{pmatrix}
        1 & 1 & 1 & 1 & -1\\
        4 & 3 & 5 & -1 & -1\\
        a & 1 & 3 & b & 1
    \end{pmatrix}\overset{r}{\sim}\begin{pmatrix}
        1 & 1 & 1 & 1 & -1\\
        0 & 1 & -1 & 5 & -3\\
        0 & 1-a & 3-a & b-a & 1+a
    \end{pmatrix}\]

    所以有$\frac{1-a}{1}=\frac{3-a}{-1}=\frac{b-a}{5}=\frac{1+a}{-3}$，得$a=2,\ b=-3$. 矩阵行变换结果为$\begin{pmatrix}
        1 & 0 & 2 & -4 & 2\\
        0 & 1 & -1 & 5 & -3\\
        0 & 0 & 0 & 0 & 0
    \end{pmatrix}$，得通解为$c_1(-2,1,1,0)^\T+c_2(4,-5,0,1)^\T+(6,-8,0,1)^\T$.
\end{enumerate}
\end{solution}

\begin{problem}
    设向量$\bm{\alpha}_1=(1,0,0,1)^\T,\ \bm{\alpha}_2=(0,1,0,2)^\T,\ \bm{\alpha}_3=(0,0,1,3)^\T,\ \bm{\alpha}_4=(1,2,3,14)^\T,\ \bm{\alpha}_5=(4,5,6,32)^\T$. 求该向量的秩，并确定一个极大线性无关组，将其余向量用该极大线性无关组线性表出.
\end{problem}
\solute1
\begin{solution}
    \[(\bm{\alpha}_1,\bm{\alpha}_2,\bm{\alpha}_3,\bm{\alpha}_4,\bm{\alpha}_5)=\begin{pmatrix}
        1 & 0 & 0 & 1 & 4\\
        0 & 1 & 0 & 2 & 5\\
        0 & 0 & 1 & 3 & 6\\
        1 & 2 & 3 & 14 & 32
    \end{pmatrix}\to
    \begin{pmatrix}
        1 & 0 & 0 & 1 & 4\\
        0 & 1 & 0 & 2 & 5\\
        0 & 0 & 1 & 3 & 6\\
        0 & 2 & 3 & 13 &28
    \end{pmatrix}\to\begin{pNiceMatrix}
        \CodeBefore
            \tikz \draw [H6,thick] (2-|1) |- (2-|2) |- (3-|2) |- (3-|3) |- (4-|3) |- (4-|6) ;
        \Body
            1 & 0 & 0 & 1 & 4\\
            0 & 1 & 0 & 2 & 5\\
            0 & 0 & 1 & 3 & 6\\
            0 & 0 & 0 & 0 & 0
        \end{pNiceMatrix}\]

        该向量组的秩为$3$，$\bm{\alpha}_1,\bm{\alpha}_2,\bm{\alpha}_3$是向量组的一个极大线性无关组，且$\bm{\alpha}_4=\bm{\alpha}_1+2\bm{\alpha}_2+3\bm{\alpha}_3,\ {\alpha}_5=4\bm{\alpha}_1+5\bm{\alpha}_2+6\bm{\alpha}_3$.
\end{solution}

\begin{problem}
    已知实二次型$f(x_1,x_2,x_3)=x_1^2+x_2^2+5x_3^2+2tx_1x_2-2x_1x_3+4x_2x_3$. 求当$t$取何值时，该二次型是正定的.
\end{problem}
\solute0
\begin{solution}*
    二次型正定的条件是
\[
    \mathrm{H}_1=1>0,  \mathrm{H}_2=\begin{vmatrix}
        1 & t\\ t & 1
    \end{vmatrix}=1-t^2>0,\ -1<t<1,\ \mathrm{H}_3=\begin{vmatrix}
        1 & t & -1\\
        t & 1 & 2\\
        -1 & 2 & 5
    \end{vmatrix}=-5t^2-4t>0,\ -\frac45<t<0\]

    所以该二次型正定的条件是$-4/5<t<0$.
\end{solution}

\begin{problem}
    求当$a,\ b$取什么值时，线性方程组
    $\begin{cases}
        x_1+x_2+x_3+x_4+x_5=1\\
        3x_1+2x_2+x_3+x_4-3x_5=a\\
        x_2+2x_3+2x_4+6x_5=3\\
        5x_1+4x_2+3x_3+3x_4-x_5=b
    \end{cases}$
    有解，并在有解的情形下求一般解.
\end{problem}
\solute1
\begin{solution}
    对线性方程组的增广矩阵进行初等行变换
    \[\mat{A}=\begin{pNiceArray}{ccccc:c}
        1 & 1 & 1 & 1 & 1 & 1\\
        3 & 2 & 1 & 1 & -3& a\\
        0 & 1 & 2 & 2 & 6 & 3\\
        5 & 4 & 3 & 3 & -1& b
    \end{pNiceArray}\to
    \begin{pNiceArray}{ccccc:c}
        1 & 1 & 1 & 1 & 1 & 1\\
        0 & -1 & -2 & -2 & -6 &  a-3\\
        0 & 1 & 2 & 2 & 6 & 3\\
        0 & -1 & -2 & -2 & -6 &  b-5
    \end{pNiceArray}\to
    \begin{pNiceArray}{ccccc:c}
        1 & 1 & 1 & 1 & 1 & 1\\
        0 & 1 & 2 & 2 & 6 & 3-a\\
        0 & 0 & 0 & 0 & 0 & a\\
        0 & 0 & 0 & 0 & 0 & b-a-2
    \end{pNiceArray}
    \]

    当$a=0,\ b=2$时，$\rank{\mat{A}}=2$，方程组有解. 
    
    取$(x_3,x_4,x_5)^\T$为自由变量，令$(x_3,x_4,x_5)^\T=(k_1,k_2,k_3)^\T$，则一般解为
    \[x=k_1(1,-2,1,0,0)^\T+k_2(1,-2,0,1,0)^\T+(5,-6,0,0,1)^\T+(-2,3,0,0,0)^\T\]
\end{solution}

\begin{problem}
    求当$x$为何值时，矩阵$\mat{A}=\begin{pmatrix}
        0 & 0 & 1\\ 1 & 1 & x\\ 1 & 0 & 0
    \end{pmatrix}$能相似对角化.
\end{problem}
\solute1
\begin{solution}*
    $\abs{\lambda\mat{E}-\mat{A}}=(\lambda-1)^2(\lambda+1)$，$\mat{A}$的特征值为$\lambda_1=\lambda_2=1,\ \lambda_3=-1$.
    \[\mat{E}-\mat{A}=\begin{pmatrix}
        1 & 0 & -1\\
        -1 & 0 & -x\\
        -1 & 0 & 1
    \end{pmatrix}\to
    \begin{pmatrix}
        1 & 0 & -1\\
        0 & 0 & x+1\\
        0 & 0 & 0
    \end{pmatrix}\]

    由$\mat{A}$能对角化，对应特征值$1$的线性无关特征向量有两个，所以$\rank{\mat{E}-\mat{A}}=1$，即$x=-1$.
\end{solution}

\begin{problem}
    设$\bm{\alpha}_1,\ \bm{\alpha}_2,\ \cdots,\ \bm{\alpha}_n$是$n$个$n$维线性无关的列向量，有$n$维列向量$\bm{\beta}$与$\bm{\alpha}_1,\ \bm{\alpha}_2,\ \cdots,\ \bm{\alpha}_n$正交. 试证明：$\bm{\beta}=0$.
\end{problem}
\solute1
\begin{solution}
方法1.
\begin{proof}
由于$\bm{\alpha}_1,\ \bm{\alpha}_2,\ \cdots,\ \bm{\alpha}_n$是$n$个$n$维线性无关的列向量，则任意一个$n$维向量均可以用$\bm{\alpha}_1,\ \bm{\alpha}_2,\ \cdots,\ \bm{\alpha}_n$线性表示. 设$\bm{\beta}=k\bm{\alpha}_1+k_2\bm{\alpha}_2+\cdots+k_n\bm{\alpha}_n$，同时与$\bm{\beta}$作内积，$\bm{\alpha}_i\bot\bm{\beta}$，所以$(\bm{\alpha}_i,\bm{\beta})=0$，即$\norm{\bm{\beta}}^2=0,\ \bm{\beta}=0$
\end{proof}
\vskip1ex\hrule\vskip1ex
方法2.
\begin{proof}
    令$\mat{A}=(\bm{\alpha}_1^\T,\bm{\alpha}_2^\T,\cdots,\bm{\alpha}_n^\T)^\T$. 由$\bm{\alpha}_i$与$\bm{\beta}$正交，所以$(\bm{\alpha}_i,\bm{\beta})=0,\ \mat{A}\bm{\beta}=\bf{0}$. 又因为$\mat{A}$由$n$个$n$维线性无关的行向量拼成，$\rank{\mat{A}}=n$，所以$\mat{A}\bm{x}=0$只有零解，$\bm{\beta}=0$.
\end{proof}
\end{solution}

\section{综合题}
\begin{problem}
    用正交线性替换将实二次型$f(x_1,x_2,x_3)=5x_1^2+5x_2^2+2x_3^3+2x_1x_2+4x_1x_3-4x_2x_3$化为标准型，并写出正交线性替换.
\end{problem}
\begin{solution}
    $\mat{A}=\begin{pmatrix}
        5 & 1 & 2\\
        1 & 5 & -2\\
        2 & -2 & 2
    \end{pmatrix},\ \lambda\mat{E}-\mat{A}=\begin{pmatrix}
        \lambda-5 & -1 & -2\\
        -1 & \lambda-5 & 2\\
        -2 & 2 & \lambda-2
    \end{pmatrix},\ \abs{\lambda\mat{E}-\mat{A}}=\lambda(\lambda-6)^2$. $\mat{A}$的特征值为$\lambda_1=0,\ \lambda_2=\lambda_3=6$.
\begin{itemize}
    \item $\lambda_1=0$，代入$\lambda_1\mat{E}-\mat{A}$，有$\mat{O}-\mat{A}=\begin{pmatrix}
        -5 & -1 & -2\\
        -1 & -5 & 2\\
        -2 & 2 & -2
    \end{pmatrix}\to
    \begin{pmatrix}
        1 & 5 & -2\\
        0 & 2 & -1\\
        0 & 0 & 0
    \end{pmatrix}$.
    
    由$(\lambda_1\mat{E}-\mat{A})\mat{X}=\mat{O}$得$\xi_1=(1,-1,-2)^\T$.
    \item $\lambda_2=\lambda_3=6$，代入$\lambda\mat{E}-\mat{A}$，有$6\mat{E}-\mat{A}=\begin{pmatrix}
        1 & -1 & -2\\
        -1 & 1 & 2\\
        -2 & 2 & 4
    \end{pmatrix}\to
    \begin{pmatrix}
        1 & -1 & -2\\
        0 & 0 & 0\\
        0 & 0 & 0
    \end{pmatrix}$.
    
    由$(6\mat{E}-\mat{A})\mat{X}=\mat{O}$得$\xi_2=(1,1,0)^\T,\ \xi_3=(2,0,1)^\T$.
    \item 进行施密特正交化 \& 单位化
    \[
    \begin{aligned}
        &\begin{cases}
            \xi_2'=\xi_2\\
            \xi_3'=\xi_3-\frac{(\xi_3,\xi_3')}{(\xi_2',\xi_2')}\cdot\xi_2=(1,-1,1)^\T
        \end{cases}\\
        &\begin{cases}
            \eta_1=\frac{\xi_1}{\norm{\xi_1}}=\frac{1}{\sqrt6}(1,-1,-1)^\T\\
            \eta_2=\frac{\xi_2}{\norm{\xi_3}}=\frac{1}{\sqrt2}(1,1,0)^\T\\
            \eta_3=\frac{\xi_2}{\norm{\xi_3}}=\frac{1}{\sqrt3}(1,-1,-1)^\T
        \end{cases}
    \end{aligned}
    \]
    \item 得正交变换
    \[\mat{Q}=\begin{pmatrix}
        \tfrac{1}{\sqrt6} & \tfrac{1}{\sqrt2} & \tfrac{1}{\sqrt3}\\
        -\tfrac{1}{\sqrt6} & \tfrac{1}{\sqrt2} & -\tfrac{1}{\sqrt3}\\
        -\tfrac{1}{\sqrt6} & 0 & \tfrac{1}{\sqrt3}
    \end{pmatrix},\ 
    \begin{pmatrix}
        x_1\\x_2\\x_3
    \end{pmatrix}=\mat{Q}
    \begin{pmatrix}
        y_1\\y_2\\y_3
    \end{pmatrix}\]
    
    得标准型为$f(x_1,x_2,x_3)=6y_2^2+6y_3^2$.
\end{itemize}
\end{solution}


\vfill
\rightline{/*Powered by \pkg{LiteSolution} \textsf\LaTeX~Template. Designed by WeYoung \href{https://ctan.org/home/xiamyphys}{Axia}.*/}